% Rozdział aktywny po restrukturyzacji DT.
\section{Rozwiązania fundamentalne operatorów różniczkowych}
\providecommand{\atlasref}[2]{#2}

\subsection{Definicja i związek ze splotem}

\begin{definicja}[Rozwiązanie fundamentalne]
\label{def:dt-fundamental-solution}
Niech $P(D)=\sum_{|\alpha|\le m}a_\alpha\partial^\alpha$
będzie operatorem różniczkowym o stałych współczynnikach
$a_\alpha\in\K$ na $\R^d$. Przez $D$ w tym zapisie
rozumiemy zwykłe pochodne cząstkowe, bez czynnika $i$.
Dystrybucję $E\in\mathcal D'(\R^d)$ nazywamy
\emph{rozwiązaniem fundamentalnym} operatora $P(D)$,
jeżeli
\[
 P(D)E=\delta_0.
\]
Równość rozumiemy w sensie pochodnych dystrybucyjnych
(Definicja~\ref{def:dist-deriv20}).
\end{definicja}

\begin{stwierdzenie}[Rozwiązanie równania z prawą stroną o zwartym nośniku]
\label{prop:dt-fundamental-convolution}
Niech $E$ będzie rozwiązaniem fundamentalnym $P(D)$.
Dla każdego $S\in\mathcal E'(\R^d)$ dystrybucja
$u=E*S$ jest rozwiązaniem równania $P(D)u=S$.
Jeżeli $S=T_f$ dla $f\in\mathcal D(\R^d)$, to $u$
jest reprezentowana przez funkcję gładką
\[
 u(x)=\langle E_y,f(x-y)\rangle,
\]
która spełnia równanie $P(D)u=f$ klasycznie.
Wszystkie dystrybucyjne rozwiązania dla ustalonego $S$
mają postać $E*S+v$, gdzie $P(D)v=0$.
W szczególności każde dwa rozwiązania fundamentalne
różnią się o rozwiązanie równania jednorodnego.
\end{stwierdzenie}
\begin{proof}
Splot jest określony, ponieważ $S$ ma zwarty nośnik
(Definicja~\ref{def:convolution21} i
Twierdzenie~\ref{thm:convolution-algebra21}).
Ze stałości współczynników i wzorów na pochodne
splotu (Twierdzenie~\ref{thm:convolution-derivative21})
otrzymujemy
\[
 P(D)(E*S)=(P(D)E)*S=\delta_0*S=S,
\]
gdzie ostatnia równość wynika z
Przykładu~\ref{ex:delta-convolution21}.
Dla $S=T_f$ Twierdzenie~\ref{thm:smoothing21}
daje wskazaną funkcję gładką. Zgodność pochodnych
klasycznych i dystrybucyjnych
(Stwierdzenie~\ref{prop:classical-weak20})
oraz różnowartościowość przyporządkowania funkcji
ciągłej odpowiadającej jej dystrybucji regularnej
dają równanie punktowe. Ostatnie twierdzenia
wynikają z liniowości $P(D)$: różnica dwóch
rozwiązań należy do jego jądra, a dodanie
elementu jądra nie zmienia prawej strony.
\end{proof}

\begin{uwaga}
Powyższe stwierdzenie nie zakłada istnienia
rozwiązania fundamentalnego dla każdego operatora.
Nie uzasadnia też splotu $E*S$ dla dowolnej
dystrybucji $S$ bez dodatkowych założeń
(por.~Uwaga~\ref{rem:convolution-limits21}).
Rozwiązanie fundamentalne nie jest automatycznie
jądrem funkcji Greena dla zagadnienia brzegowego:
nie narzucono tutaj żadnych warunków brzegowych.
\end{uwaga}

\subsection{Warunek istnienia rozwiązania o zwartym nośniku}

\begin{twierdzenie}[Rozwiązanie fundamentalne o zwartym nośniku]
\label{thm:dt-fundamental-compact}
Operator różniczkowy o stałych współczynnikach $P(D)$
na $\R^d$ ma rozwiązanie fundamentalne o zwartym
nośniku wtedy i tylko wtedy, gdy jest mnożeniem
przez niezerową stałą $a$. Wówczas jedynym
rozwiązaniem fundamentalnym jest $a^{-1}\delta_0$.
\end{twierdzenie}
\begin{proof}
Przypadek niezerowej stałej jest natychmiastowy,
a operator zerowy nie może mieć rozwiązania fundamentalnego.
Załóżmy zatem, że wielomian symbolu
$P(z)=\sum_\alpha a_\alpha z^\alpha$ ma stopień $m\ge1$.
Jego część jednorodna najwyższego stopnia $P_m$
nie jest wielomianem zerowym. Istnieje $v\in\R^d$
z $P_m(v)\ne0$: wielomian znikający na całym
$\R^d$ jest zerowy, co wynika przez indukcję
względem liczby zmiennych z jednowymiarowego
twierdzenia o liczbie pierwiastków wielomianu.
Wielomian $t\mapsto P(tv)$ ma stopień $m$, więc
z zasadniczego twierdzenia algebry istnieje
$t_0\in\C$ takie, że $P(t_0v)=0$.
Oznaczmy $z_0=t_0v$.

Jeżeli $E$ ma zwarty nośnik i $P(D)E=\delta_0$,
to równość zachodzi również w $\mathcal E'(\R^d)$
na mocy Twierdzenia~\ref{thm:compact-support-E25}.
Nad $\R$ przedłużamy działanie dystrybucji
zespolonoliniowo na funkcje zespolone.
Możemy więc zastosować obie strony do gładkiej
funkcji $h(x)=\exp(-z_0\cdot x)$, która nie musi
mieć zwartego nośnika. Z definicji pochodnej
i przedłużenia działania na funkcje gładkie
(Lemat~\ref{lem:compact-action21}) otrzymujemy
\[
 1=h(0)=\langle P(D)E,h\rangle
 =\langle E,P(-D)h\rangle
 =P(z_0)\langle E,h\rangle=0.
\]
Ostatnią równość poprzedza obliczenie
$\partial^\alpha h=(-z_0)^\alpha h$.
Sprzeczność dowodzi tezy.
\end{proof}

\subsection{Równanie symbolu w przestrzeni dystrybucji temperowanych}

\begin{stwierdzenie}[Charakterystyka przez transformację Fouriera]
\label{prop:dt-fundamental-symbol}
Przy normalizacji $c_d=(2\pi)^{-d/2}$ dystrybucja
$E\in\mathcal S'(\R^d)$ jest rozwiązaniem fundamentalnym
$P(D)$ wtedy i tylko wtedy, gdy
\[
 P(i\xi)\widehat E=c_d T_1.
\]
\end{stwierdzenie}
\begin{proof}
Wzory na transformację pochodnych dystrybucyjnych
i delty Diraca (Przykład~\ref{ex:fourier-delta22})
dają $\mathcal F(P(D)E)=P(i\xi)\widehat E$
oraz $\mathcal F\delta_0=c_d T_1$.
Implikacja odwrotna wynika z różnowartościowości
transformacji Fouriera na $\mathcal S'$
(Twierdzenie~\ref{thm:fourier-auto-tempered22}).
Równość w $\mathcal D'$ jest równoważna równości
w $\mathcal S'$ dla tych dystrybucji, przez
różnowartościowość ograniczenia do $\mathcal D$
(Stwierdzenie~\ref{prop:schwartz-embeddings22}).
\end{proof}

\begin{uwaga}
Równanie symbolu jest równaniem mnożenia dystrybucji
przez wielomian, a nie zwykłym dzieleniem funkcji.
Gdy symbol ma zera, funkcja $1/P(i\xi)$ może
nie być lokalnie całkowalna. Charakterystyka
nie stanowi sama w sobie dowodu istnienia
rozwiązania fundamentalnego dla każdego niezerowego
operatora.
\end{uwaga}

\subsection{Rozwiązanie fundamentalne operatora Laplace'a}
\begin{twierdzenie}[Rozwiązanie fundamentalne operatora $-\Delta$]
\label{thm:dt-laplace}
Niech $d\ge1$ i niech $\omega_d$ oznacza pole powierzchni
sfery jednostkowej $S^{d-1}\subset\R^d$ dla $d\ge2$.
Funkcja lokalnie całkowalna
\begin{equation}\label{eq:dt-laplace-fundamental}
 \Phi_d(x)=\begin{cases}
 -|x|/2,&d=1,\\
 (2\pi)^{-1}\log(1/|x|),&d=2,\\
 ((d-2)\omega_d)^{-1}|x|^{2-d},&d\ge3,
 \end{cases}
\end{equation}
spełnia $-\Delta T_{\Phi_d}=\delta_0$ na $\R^d$.
\end{twierdzenie}
\begin{proof}
Wartość $\Phi_d(0)$ w wymiarach $d\ge2$ ustalamy
dowolnie, na przykład równą $0$.
W wymiarze jeden $\Phi_1$ jest lokalnie całkowalna
i jej pierwsza pochodna dystrybucyjna jest funkcją
$-\tfrac12\operatorname{sgn}x$, ponieważ funkcja jest ciągła
i kawałkami liniowa. Skok tej pochodnej w zerze wynosi $-1$.
Z Twierdzenia~\ref{thm:jump20} wynika zatem
$\partial_x^2T_{\Phi_1}=-\delta_0$.

Dla $d\ge3$ całkowanie w zmiennych biegunowych daje
dla $r>0$
\[
 \int_{B_r(0)}|\Phi_d(x)|\,\mathrm d\lambda_d(x)
 =\frac1{d-2}\int_0^r s\,\mathrm d\lambda_1(s)
 =\frac{r^2}{2(d-2)}.
\]
Dla $d=2$ mamy $\omega_2=2\pi$ i dla $0<r<1$
\[
 \begin{aligned}
 \int_{B_r(0)}|\Phi_2(x)|\,\mathrm d\lambda_2(x)
 &=\int_0^r s\log(1/s)\,\mathrm d\lambda_1(s)\\
 &=\frac{r^2}{2}\log(1/r)+\frac{r^2}{4}.
 \end{aligned}
\]
Granica dolnego wyrazu pierwotnej jest zerowa,
gdyż $s^2\log(1/s)\to0$ dla $s\downarrow0$.
W obu przypadkach całka dąży do zera przy $r\downarrow0$.
Z gładkości poza zerem wynika więc lokalna
całkowalność na całej $\R^d$.

Niech $F(r)$ oznacza profil radialny $\Phi_d$ dla $r>0$.
Bezpośrednie różniczkowanie obu wzorów daje
\[
 F'(r)=-\frac1{\omega_dr^{d-1}},\qquad
 F''(r)=\frac{d-1}{\omega_dr^d}.
\]
Dla $x\ne0$, $r=|x|$,
$\partial_jr=x_j/r$ oraz
$\partial_j^2r=1/r-x_j^2/r^3$.
Z reguły łańcuchowej otrzymujemy
\[
 \partial_j^2\Phi_d(x)
 =F''(r)\frac{x_j^2}{r^2}
 +F'(r)\left(\frac1r-\frac{x_j^2}{r^3}\right).
\]
Sumując po $j$ i korzystając z $\sum_j x_j^2=r^2$,
dostajemy
\[
 \Delta\Phi_d(x)
 =F''(r)+\frac{d-1}{r}F'(r)=0.
\]
Dla $\varphi\in\mathcal D(\R^d)$ wybieramy $R>0$
tak, aby $\supp\varphi\Subset B_R(0)$.
Dla $0<\varepsilon<R$ stosujemy drugą tożsamość
Greena (wynikającą z twierdzenia o dywergencji
dla pola $\Phi_d\nabla\varphi-\varphi\nabla\Phi_d$) na pierścieniu
$A_{\varepsilon,R}=B_R(0)\setminus\overline{B_\varepsilon(0)}$
(do części rzeczywistej i urojonej testu osobno).
Ponieważ $\Delta\Phi_d=0$ na pierścieniu,
otrzymujemy
\[
 -\int_{A_{\varepsilon,R}}\Phi_d\Delta\varphi
 \,\mathrm d\lambda_d
 =\int_{\partial A_{\varepsilon,R}}
 (\varphi\partial_\nu\Phi_d-\Phi_d\partial_\nu\varphi)
 \,\mathrm d\mathcal H^{d-1}.
\]
Na zewnętrznej sferze znikają test i jego pochodne.
Na wewnętrznej sferze $\nu=-x/|x|$, zatem
$\partial_\nu\Phi_d=-F'(\varepsilon)
=1/(\omega_d\varepsilon^{d-1})$.
W konsekwencji
\[
 -\int_{A_{\varepsilon,R}}\Phi_d(x)\Delta\varphi(x)
 \,\mathrm d\lambda_d(x)
 =\int_{\partial B_\varepsilon}\varphi(x)
 \frac{1}{\omega_d\varepsilon^{d-1}}\,\mathrm d\mathcal H^{d-1}(x)
 -\int_{\partial B_\varepsilon}\Phi_d(x)
 \partial_\nu\varphi(x)\,\mathrm d\mathcal H^{d-1}(x).
\]
Miara powierzchniowa sfery o promieniu $\varepsilon$
wynosi $\omega_d\varepsilon^{d-1}$, przez jednorodność
miary Hausdorffa (por. dowód Lematu~\ref{lem:dt-fourier-sphere-kernel}). Po odjęciu
$\varphi(0)$ od pierwszej całki moduł różnicy
jest więc co najwyżej
$\sup_{|x|=\varepsilon}|\varphi(x)-\varphi(0)|$,
co dąży do zera z ciągłości testu.
Moduł drugiej całki nie przekracza
\[
 \|\nabla\varphi\|_\infty\,
 |F(\varepsilon)|\,\omega_d\varepsilon^{d-1}.
\]
Dla $d\ge3$ iloczyn bez normy gradientu wynosi
$\varepsilon/(d-2)$, dla $d=2$ i $\varepsilon<1$
wynosi $\varepsilon\log(1/\varepsilon)$.
Obie wielkości dążą do zera.
Ponieważ sfery mają miarę Lebesgue'a zero
i test znika poza $B_R(0)$, w całce po $\R^d$
pozostaje uwzględnić tylko wkład z $B_\varepsilon(0)$.
Jego moduł jest ograniczony przez
$\|\Delta\varphi\|_\infty
\int_{B_\varepsilon(0)}|\Phi_d(x)|\,\mathrm d\lambda_d(x)$,
co dąży do zera z wcześniejszych obliczeń.
Ze znaku $(-1)^2=1$ przy drugiej pochodnej
dystrybucyjnej otrzymujemy
\[
 \langle-\Delta T_{\Phi_d},\varphi\rangle
 =-\int_{\R^d}\Phi_d(x)\Delta\varphi(x)
 \,\mathrm d\lambda_d(x)=\varphi(0).
\]
Jest to żądana równość.
\end{proof}

\begin{uwaga}[Powiązanie z równaniami różniczkowymi cząstkowymi]
\noindent W skrypcie RRC rozwijamy zastosowania do
zagadnień brzegowych i ewolucyjnych:
\atlasref{PDE-03::sec:pde3-green}{funkcje Greena i tożsamości Greena},
\atlasref{PDE-07::thm:pde7-heat}{jądro cieplne jako rozwiązanie fundamentalne}
oraz
\atlasref{PDE-07::thm:pde7-weyl}{lemat Weyla}.
Tutaj koncentrujemy się na rachunku dystrybucji,
splotach, symbolach i ograniczeniach wynikających
z nośnika.
\end{uwaga}
